\documentclass[UTF8,fontset=fandol,10pt]{ctexart}
\usepackage[a4paper,margin=22mm]{geometry}
\usepackage{amsmath,amssymb,booktabs,tabularx,array,multirow,graphicx,cite,xurl,placeins}
\usepackage{longtable,xltabular,needspace}
\setmainfont{TeX Gyre Termes}
\emergencystretch=2em
\pagestyle{plain}
\linespread{1.15}
\renewcommand{\thetable}{S\arabic{table}}
\title{RGSP 补充材料\\协议、逐种子结果与历史记录}
\author{}
\date{2026 年 10 月 9 日修订}
\begin{document}
\maketitle
\section{阅读范围与证据等级}
本材料与正文分工：正文呈现计算机制与可解释的配置比较，本材料保存逐 run 结果、控制定义和历史记录。历史表格的原始数值未因结果不利而删除，也没有升级为同协议排行榜。源自不同论文的外部数值属于旧稿收集记录，完整逐行主文献数值核验尚未完成；具体来源、已核对的书目信息和缺口见随包 \texttt{SOURCE\_REGISTRY.md}。仅补上文献名称或 DOI，不表示已复现该方法或核验了其全部数值。

本版依据本地最新双栏正文包及 10 月 8 日两份图件修订；原 26 页单栏 PDF 与当前可编辑正文范围不完全相同，本材料不声称恢复该旧 PDF 的全部附录。原始 ZIP、旧稿及实验记录在原位置保留，新包的历史源副本另存于 \texttt{archive/20261008}，不作为当前结论使用。

\section{控制实验协议}
SSDD 历史 15-run 为 742/186/232 轨道，存在关闭参数漂移及内容重叠限制；新的 730/186/232 轨道保留官方 test，按内容组划分开发集，采用运行时关闭条件与初始化核验。字节去重不能替代近重复、场景关联和标签一致性审计。新轨道完整结果在本次回收范围内尚未验收，不以部分 run 或旧成绩填充。

\begin{table}[!htbp]
\centering\small
\caption{新 SSDD 的 12 配置定义。R 为残差修正，E 为趋势能量调制；U、$\Delta$ 为结构引导作用路径。}
\begin{tabularx}{\linewidth}{@{}lccccX@{}}
\toprule
配置 & R&E&$I\to U$&$I\to\Delta$&说明\\
\midrule
Original & --&--&--&--&原始四向 SS2D 系统参考\\
MatchedIdentity&0&0&0&0&相同替换结构，输入倍率 1\\
Matched075&0&0&0&0&相同替换结构，输入倍率 0.75\\
Residual075&1&0&0&0&保持常数缩放背景，只开残差修正\\
EnergyOnly&0&1&0&0&只开启趋势能量调制\\
LocalEnergy&1&1&0&0&残差与能量组合\\
LocalEnergy+Input&1&1&1&0&仅增加输入结构引导\\
LocalEnergy+Delta&1&1&0&1&仅增加步长结构引导\\
Full&1&1&1&1&完整两条结构引导路径\\
FixedHalf&1&1&常数&常数&$I=0.5$，保留 0.05 步长 logits 偏置\\
GuidanceOnly&0&0&1&1&严格关闭 R/E 的引导控制\\
Global&1&1&均值&均值&可训练引导空间均值后广播\\
\bottomrule
\end{tabularx}
\end{table}
\FloatBarrier

关闭参数通过不参与优化的 buffer 固定，保持统一模块构造顺序和共同张量初始化；不是简单设置 \texttt{requires\_grad=False}。关闭分支保留的 buffer 不计为可训练参数，故不声称所有变体具有完全相同的可训练参数数量。残差/能量四格及输入/步长四格均按种子计算差分中的差分，不能用 Full 减某单一变体替代严格交互。

当前检测 Full 的计算路径对应 PFIM、DeSpeckleSSM 和 DS\_VSSBlock：PFIM 对应 LSGM；DeSpeckleSSM 是含 LRC、输入/步长调制及输出融合的整个状态建模核心，而非 LRC 的同义词；DS\_VSSBlock 包含核心之外的归一化、残差和 FFN。论文符号 $a_1,a_2$ 对应块级通道缩放，不能与扫描输入因子 $\gamma$ 混淆。正文图 2(c) 展示的是核心，不是整个替换块。

SSDD 使用 150e/patience30；SADD 使用 200e/patience40；其余增强、AdamW、线性学习率、FP32 与选模定义见正文。验证采用 $0.1\mathrm{AP}_{50}+0.9\mathrm{AP}_{50:95}$ 选模；训练验证 NMS IoU=0.7，最终独立测试 conf=0.001、NMS IoU=0.6。这两阶段的后处理设置不完全相同。

\section{SADD：全部种子与恢复边界}
固定阳性子集 795/44/44，四配置各三个训练种子。下表由回收 JSON 中的比例值乘以 100 生成，主文采用均值及样本标准差。完整比例值、Precision/Recall 和权重 SHA256 见 \texttt{sadd\_results\_individual.csv}，源 JSON 随包保留。
\begin{table}[!htbp]
\centering\small
\caption{SADD 新轨道的 12 项结果（\%），不与后文历史 SADD 表混算。}
\begin{tabular}{lrrrr}
\toprule
Configuration & Seed & AP$_{50}$ & AP$_{75}$ & AP$_{50:95}$\\
\midrule
Original & 0 & 98.3691 & 53.5163 & 53.7117 \\
Matched075 & 0 & 97.5739 & 49.5962 & 52.0914 \\
LocalEnergy & 0 & 97.5941 & 48.1029 & 51.9479 \\
Full & 0 & 97.0944 & 43.5558 & 49.0030 \\
Original & 1 & 96.8641 & 56.8648 & 53.9804 \\
Matched075 & 1 & 98.5252 & 51.6064 & 53.0962 \\
LocalEnergy & 1 & 96.6499 & 47.3977 & 52.3537 \\
Full & 1 & 98.0804 & 49.2254 & 53.1989 \\
Original & 2 & 96.7214 & 48.0764 & 52.3480 \\
Matched075 & 2 & 96.0201 & 50.8144 & 51.3444 \\
LocalEnergy & 2 & 98.4768 & 50.7618 & 52.8014 \\
Full & 2 & 98.0135 & 54.5365 & 53.7361 \\
\bottomrule
\end{tabular}
\end{table}
\FloatBarrier
Full 相对 Original 的配对差为 $-4.7087/-0.7816/+1.3881$ 个百分点。仅一个种子为正，平均 $-1.3674$，不能宣称稳定改善。Original 经 EMA/FP16 优化器断点恢复，非逐位无中断等价；seed0/1 恢复前部分控制审计记录缺失。阳性筛选排除了空场景，测试集仅 44 张，且没有完成场景独立性证明。完成训练及测试不消除这些限制。

\section{SSDD：历史标签纠正与全部记录}
旧表的 DS-only 在本稿称 LocalEnergy，包含残差修正和趋势能量调制；PFIM-only 改称 Legacy-PFIM，因为三个旧检查点中的 $\lambda_s$、$\alpha$ 非零，不满足关闭条件。原 15-run 全指标和权重摘要保存在 \path{results_individual.csv}，数值未改动。所有种子一视同仁地保留，不使用“outlier”排除 seed2。Full 减 LocalEnergy 只是边际差，不命名为严格交互。

\begin{table}[!htbp]
\centering\small
\caption{历史 SSDD 相对 Original 的 AP$_{50:95}$ 配对差（百分点）。Legacy-PFIM 不是严格仅引导控制。}
\begin{tabular}{lrrrr}
\toprule
配置 & Seed 0 & Seed 1 & Seed 2 & Mean$\pm$Std\\
\midrule
LocalEnergy & +2.90 & +5.26 & -2.30 & $+1.95\pm3.87$\\
Legacy-PFIM & +2.13 & +3.16 & -2.00 & $+1.09\pm2.73$\\
Full & -0.35 & +5.93 & -2.94 & $+0.88\pm4.57$\\
\bottomrule
\end{tabular}
\end{table}
\FloatBarrier

\section{保留的历史跨任务与文献结果}
以下表格继承最新双栏稿的数值和批次标签。历史 SADD、SSDD 高分轨道及跨任务结果均不冒充本版新控制实验；文献与项目结果仅按来源分区，不进行统一排名。未完成逐行来源核验的外部行只作为待核验记录保留，不能直接移回投稿主表。
% Retained numerical records from the latest local two-column ZIP.
% Not a new benchmark or a claim of primary-source numerical verification.
\FloatBarrier
\subsection{图像分类}
\label{sec:classification_results}

在 SAR-ACD 和 SAMPLE 数据集上评估 RGSP 的图像分类表现，
并汇总相关方法的文献报告结果。
采用分类准确率（Accuracy）和宏平均 F1
（Macro-F1）作为评价指标。
对于文献中按正确分类样本数占测试样本总数计算的
总体准确率（Overall Accuracy, OA），
本文统一记为 Accuracy。
各表按结果来源与实验设置分组；
不同类别划分、训练数据比例和初始化设置的结果
不进行统一排名。
本项目结果均为单种子记录。

\subsubsection{SAR-ACD}
\label{sec:sar_acd_results}

表~\ref{tab:sar_acd_comparison}
列出了 RGSP 与相关方法在 SAR-ACD 上的分类结果，
比较方法包括经典卷积网络、Transformer
以及面向 SAR 图像的分类与表征学习方法。
分类准确率统一记为 Accuracy，
文献中定义相同的总体准确率 OA 使用相同记法。

RGSP 采用 VMamba-Tiny 预训练权重初始化骨干，
新增模块预热 50 轮后，
Accuracy 和 Macro-F1 分别达到
96.850\% 和 96.830\%。
对应基线的两项指标分别为
89.764\% 和 89.585\%。
由于两种配置的模块预热日程不同，
该差异反映 RGSP 与相应训练策略的组合表现。
外部方法采用各自文献中的实验设置，
其结果作为参考列示。

\Needspace{10\baselineskip}
\begingroup
\small
\setlength{\tabcolsep}{5pt}
\setlength{\LTcapwidth}{\linewidth}
\begin{longtable}{@{}lrr@{}}
\caption{SAR-ACD 分类结果（\%）。
外部方法采用文献报告值，
其数据划分与训练设置未全部统一。
粗体仅标识本文方法，不表示跨协议最优。}\label{tab:sar_acd_comparison}\\
\toprule
Method & Accuracy & Macro-F1 \\
\midrule
\endfirsthead
\multicolumn{3}{c}{表~\thetable（续）}\\
\toprule
Method & Accuracy & Macro-F1 \\
\midrule
\endhead
\midrule
\multicolumn{3}{r}{接下页}\\
\endfoot
\bottomrule
\endlastfoot
AlexNet$^{a}$            & 91.65 & -- \\
GoogLeNet$^{a}$          & 93.61 & -- \\
VGGNet$^{a}$             & 92.30 & -- \\
ResNet$^{a}$             & 92.79 & -- \\
ResNet-50$^{b}$          & 84.08 & -- \\
ResNet-101$^{b}$         & 80.70 & -- \\
A-ConvNet$^{a}$          & 90.12 & -- \\
Lm-BN-CNN$^{a}$          & 87.34 & -- \\
CBAM Net$^{a}$           & 92.66 & -- \\
VIB Net$^{a}$            & 93.50 & -- \\
Swin-T$^{b}$             & 85.74 & -- \\
Swin-B$^{b}$             & 86.24 & -- \\
ViT$^{b}$                & 93.20 & -- \\
SUMMIT$^{b}$             & 94.86 & -- \\
Bottom-Up Scattering Network$^{a}$
                        & 95.43 & -- \\
SARMAE (ViT-B)$^{c}$     & 95.06 & -- \\
SARMAE (ViT-L)$^{c}$     & 95.63 & -- \\
MDB-DPL$^{d}$           & $98.85 \pm 0.17$ & -- \\
\midrule
VMamba-Tiny baseline    & 89.764 & 89.585 \\
\textbf{RGSP (Ours)}     & 96.850 & 96.830 \\
\end{longtable}
\noindent
\smallskip
\begin{minipage}{0.98\linewidth}
\footnotesize
\textbf{表~\ref{tab:sar_acd_comparison} 注：}
$^{a}$结果取自 Bottom-Up 论文的 Table 7，
保留原表中的网络名称，其中 ResNet 的具体版本未在表中注明；
$^{b}$结果取自 SUMMIT 论文的 Table 5，
采用其完整训练集设置；
$^{c}$采用 30\% 标注训练数据；
$^{d}$采用六类、8:2 划分，报告五次实验的均值与标准差。
本文采用五类设置，报告单种子结果；
RGSP 使用额外 50 轮模块预热，基线无该额外日程。
``--'' 表示未收集到对应的 Macro-F1。
\end{minipage}
\endgroup
\FloatBarrier
\subsubsection{SAMPLE}
\label{sec:sample_results}

表~\ref{tab:sample_comparison}
汇总 SAMPLE 上的分类结果。
外部方法包括 ResNet-50、A-ConvNet、
CORAL、DANN、ERM、VREx、HDANet、SSR 和 MLDA。
这些结果统一取自 MLDA 论文的比较实验，
采用合成数据训练、实测数据测试的设置。
其数值表示该论文报告的结果，
并非本文对上述方法的重新实现。

本项目采用当前闭集分类划分和从头训练设置。
RGSP 的 Accuracy 和 Macro-F1
分别为 97.037\% 和 96.441\%，
对应基线分别为 98.519\% 和 98.454\%。
因此，在该设置下，
RGSP 的两项指标分别下降
1.482 和 2.013 个百分点，
尚未表现出分类性能增益。
由于当前项目划分尚未与文献中的
合成训练、实测测试协议对齐，
两组结果分别列示，不计算跨组性能增益。

\Needspace{10\baselineskip}
\begingroup
\small
\setlength{\tabcolsep}{5pt}
\setlength{\LTcapwidth}{\linewidth}
\begin{longtable}{@{}lrr@{}}
\caption{SAMPLE 分类结果（\%）。
外部结果来自 MLDA 论文
（Remote Sensing, 2025）的 Table 2，
采用合成数据训练、实测数据测试。
本项目结果采用当前闭集划分，为单种子记录。
粗体仅标识本文方法。}\label{tab:sample_comparison}\\
\toprule
Method & Accuracy & Macro-F1 \\
\midrule
\endfirsthead
\multicolumn{3}{c}{表~\thetable（续）}\\
\toprule
Method & Accuracy & Macro-F1 \\
\midrule
\endhead
\midrule
\multicolumn{3}{r}{接下页}\\
\endfoot
\bottomrule
\endlastfoot
\multicolumn{3}{l}{
\textit{Literature results: synthetic-to-measured}} \\
ResNet-50
& $48.37 \pm 3.73$ & $45.54 \pm 3.75$ \\
A-ConvNet
& $52.36 \pm 3.97$ & $48.47 \pm 3.66$ \\
CORAL
& $87.76 \pm 2.58$ & $86.92 \pm 2.67$ \\
DANN
& $90.42 \pm 2.06$ & $89.87 \pm 2.14$ \\
ERM
& $91.78 \pm 1.86$ & $90.97 \pm 1.72$ \\
VREx
& $92.64 \pm 2.09$ & $92.67 \pm 1.89$ \\
HDANet
& $86.21 \pm 2.61$ & $85.79 \pm 2.47$ \\
SSR
& $94.69 \pm 1.19$ & $94.65 \pm 1.21$ \\
MLDA
& $96.76 \pm 1.34$ & $96.70 \pm 1.29$ \\

\midrule
\multicolumn{3}{l}{
\textit{Our experiments: current closed-set split}} \\
Classification baseline
& 98.519 & 98.454 \\
\textbf{RGSP (Ours)}
& 97.037 & 96.441 \\
\end{longtable}
\endgroup
\FloatBarrier
\FloatBarrier
\subsection{目标检测}
\label{sec:detection_results}

在 SSDD、SAR-ShipDataset、SADD 和 AIR-SARShip-2.0
四个数据集上评估 RGSP 的目标检测性能，
并与通用目标检测器及 SAR 专用检测方法进行比较。
评价指标包括 AP$_{50}$、AP$_{75}$ 和 AP$_{50:95}$。
外部方法采用文献报告结果，
本文方法采用已有实验记录（非本版新 SADD 三种子结果）；
不同数据划分、初始化方式和训练预算的差异在表下注明。
所有结果均以百分数表示，
``--'' 表示未报告或尚未确认的指标。

\subsubsection{SSDD}
\label{sec:ssdd_results}

表~\ref{tab:ssdd_comparison}
汇总 SSDD 上的检测结果。
RGSP 在已有的 300 轮训练配置下，
AP$_{50}$、AP$_{75}$ 和 AP$_{50:95}$
分别达到 $98.549\%$、$85.552\%$ 和 $72.261\%$。
该结果来自官方测试集评估，
与主文中的历史三种子观察分别报告。

\Needspace{10\baselineskip}
\begingroup
\small
\setlength{\tabcolsep}{5pt}
\setlength{\LTcapwidth}{\linewidth}
\begin{longtable}{@{}lrrr@{}}
\caption{SSDD 目标检测结果（\%）。
外部方法为文献报告值；
本文列示已有官方测试集结果。
粗体方法名仅标识本文方法。}\label{tab:ssdd_comparison}\\
\toprule
Method & AP$_{50}$ & AP$_{75}$ & AP$_{50:95}$ \\
\midrule
\endfirsthead
\multicolumn{4}{c}{表~\thetable（续）}\\
\toprule
Method & AP$_{50}$ & AP$_{75}$ & AP$_{50:95}$ \\
\midrule
\endhead
\midrule
\multicolumn{4}{r}{接下页}\\
\endfoot
\bottomrule
\endlastfoot
SAR-ShipNet       & 89.10 & --    & --    \\
DCTC              & 93.90 & --    & --    \\
GEFCSI-Net        & 97.10 & --    & --    \\
SARNas            & 98.50 & --    & --    \\
MGA-Net           & 98.60 & --    & --    \\
Faster R-CNN      & 96.40 & 86.10 & 71.30 \\
Cascade R-CNN     & 96.40 & 87.40 & 72.70 \\
RetinaNet         & 95.00 & --    & --    \\
CenterNet         & 95.10 & --    & --    \\
FCOS              & 95.30 & --    & --    \\
ConsistencyDet    & 95.60 & --    & --    \\
YOLOv3            & 96.20 & --    & --    \\
YOLOX             & 96.20 & --    & --    \\
DiffusionDet      & 96.30 & --    & --    \\
YOLOv5n           & 98.00 & --    & --    \\
YOLOv5s           & 98.30 & --    & --    \\
YOLOv8n           & 97.10 & 91.40 & 74.00 \\
YOLOv9-tiny       & 97.30 & 88.90 & 72.20 \\
YOLOv10n          & 97.30 & 88.50 & 72.80 \\
DiffDet4SAR       & 96.90 & --    & --    \\
YOLO-SRBD         & 98.30 & --    & 74.50 \\
LRTransDet        & 97.80 & --    & --    \\
SFS-CNet          & 99.60 & --    & --    \\
MaDiNet           & 99.00 & --    & --    \\
CCAI-YOLO         & 98.80 & --    & 74.90 \\
ARGO-Net          & 99.50 & --    & 79.30 \\
\midrule
Mamba-YOLO-T      & 98.605 & 83.613 & 70.348 \\
\textbf{RGSP (Ours)}
                  & 98.549 & 85.552 & 72.261 \\
\end{longtable}
\noindent
\smallskip
\begin{minipage}{\linewidth}
\footnotesize
\textbf{表~\ref{tab:ssdd_comparison} 注：}
外部结果来自不同文献的比较记录。
对于存在多个报告值的方法，
本表保留原资料中同时提供更多 AP 指标的条目，
不将不同来源的指标拼接为同一行。
Mamba-YOLO-T 为历史官方测试集记录，
RGSP 为 300 轮训练配置；
二者的训练预算与初始化尚未完全对齐，
因此不据此计算受控模块增益。
\end{minipage}
\endgroup
\FloatBarrier


\subsubsection{SAR-ShipDataset}
\label{sec:sar_ship_results}

表~\ref{tab:sar_ship_comparison}
列出 SAR-ShipDataset 上的检测结果。
已有 RGSP 配置的 AP$_{50:95}$ 为 $57.818\%$。
采用较大模型规模时，
RGSP-S 的 AP$_{50}$ 为 $95.609\%$，
但其 AP$_{50:95}$ 为 $53.177\%$。
两组结果分别保留完整指标，
不将不同配置的最优指标合并为同一结果。

\Needspace{10\baselineskip}
\begingroup
\small
\setlength{\tabcolsep}{5pt}
\setlength{\LTcapwidth}{\linewidth}
\begin{longtable}{@{}lrrr@{}}
\caption{SAR-ShipDataset 目标检测结果（\%）。
本文同时列出 RGSP 与 RGSP-S 的已有结果，
各行对应一个独立配置。}\label{tab:sar_ship_comparison}\\
\toprule
Method & AP$_{50}$ & AP$_{75}$ & AP$_{50:95}$ \\
\midrule
\endfirsthead
\multicolumn{4}{c}{表~\thetable（续）}\\
\toprule
Method & AP$_{50}$ & AP$_{75}$ & AP$_{50:95}$ \\
\midrule
\endhead
\midrule
\multicolumn{4}{r}{接下页}\\
\endfoot
\bottomrule
\endlastfoot
OE-YOLO           & 86.00 & -- & --    \\
SAR-ShipNet       & 90.20 & -- & --    \\
ShipDetector      & 91.20 & -- & --    \\
Improved FCOS     & 94.10 & -- & --    \\
GEFCSI-Net        & 94.30 & -- & --    \\
CMFT              & 94.50 & -- & --    \\
YOLOv4            & 77.90 & -- & --    \\
YOLOv7            & 86.10 & -- & --    \\
YOLOv3            & 88.80 & -- & --    \\
YOLOv5n           & 89.30 & -- & --    \\
CenterNet         & 92.10 & -- & --    \\
FCOS              & 92.80 & -- & --    \\
Cascade R-CNN     & 92.70 & -- & --    \\
Faster R-CNN      & 93.10 & -- & --    \\
RetinaNet         & 93.40 & -- & --    \\
ConsistencyDet    & 94.20 & -- & --    \\
DiffusionDet      & 95.00 & -- & --    \\
LRTransDet        & 95.10 & -- & --    \\
SRDet             & 95.10 & -- & 65.90 \\
DiffDet4SAR       & 95.10 & -- & --    \\
YOLOSAR-Lite      & 96.06 & -- & --    \\
MaDiNet           & 97.50 & -- & --    \\
\midrule
Mamba-YOLO-T      & 94.616 & -- & 49.568 \\
\textbf{RGSP (Ours)}
                  & 94.844 & -- & 57.818 \\
\textbf{RGSP-S (Ours)}
                  & 95.609 & -- & 53.177 \\
\end{longtable}
\noindent
\smallskip
\begin{minipage}{\linewidth}
\footnotesize
\textbf{表~\ref{tab:sar_ship_comparison} 注：}
本文结果为历史单次实验记录，
初始化、模型规模及训练预算未全部统一。
RGSP-S 采用标准推理，不包含测试时增强。
上述结果用于展示已有配置的性能，
不作为严格受控消融结论，
也不跨协议标注最优值。
\end{minipage}
\endgroup
\FloatBarrier


\subsubsection{SADD}
\label{sec:sadd_results}

表~\ref{tab:sadd_comparison}
汇总 SADD 上的检测结果。
在已有从头训练实验中，
RGSP 的 AP$_{50:95}$ 为 $54.597\%$，
低于对应基线的 $55.018\%$，
当前结果未显示该设置下的性能增益。
此外，表中保留部分权重迁移配置的历史结果，
用于区分初始化方式带来的实验差异。

\Needspace{10\baselineskip}
\begingroup
\small
\setlength{\tabcolsep}{5pt}
\setlength{\LTcapwidth}{\linewidth}
\begin{longtable}{@{}lrrr@{}}
\caption{SADD 目标检测结果（\%）。
本文从头训练与部分权重迁移配置分别列示。}\label{tab:sadd_comparison}\\
\toprule
Method & AP$_{50}$ & AP$_{75}$ & AP$_{50:95}$ \\
\midrule
\endfirsthead
\multicolumn{4}{c}{表~\thetable（续）}\\
\toprule
Method & AP$_{50}$ & AP$_{75}$ & AP$_{50:95}$ \\
\midrule
\endhead
\midrule
\multicolumn{4}{r}{接下页}\\
\endfoot
\bottomrule
\endlastfoot
YOLOv4            & 87.80 & --    & 54.00 \\
YOLOR-P6          & 89.20 & --    & 53.50 \\
YOLOv7            & 89.80 & --    & 55.20 \\
ResNet50-CSP      & 91.50 & --    & 55.90 \\
ResNeXt50-CSP     & 91.30 & --    & 55.40 \\
YOLO-SAD          & 91.90 & --    & 57.10 \\
ViTDet            & 93.00 & 51.00 & 50.20 \\
SARATR-X          & 94.30 & 50.70 & 50.40 \\
SARDet-MIM        & 95.90 & 59.60 & 55.30 \\
\midrule
Mamba-YOLO-T (scratch)
                  & 98.335 & 57.030 & 55.018 \\
\textbf{RGSP (scratch)}
                  & 97.760 & 56.985 & 54.597 \\
\textbf{RGSP (partial transfer)}$^{\dagger}$
                  & 98.297 & 82.999 & 72.113 \\
\end{longtable}
\noindent
\smallskip
\begin{minipage}{\linewidth}
\footnotesize
\textbf{表~\ref{tab:sadd_comparison} 注：}
Scratch 表示从头训练。
$^{\dagger}$部分权重迁移结果来自历史 300 轮记录，
其权重来源与训练协议仍需完成审计，
暂作为候选结果保留，
不用于宣称相对外部方法的性能优势。
外部方法的预训练与数据划分设置并不完全一致。
\end{minipage}
\endgroup
\FloatBarrier


\subsubsection{AIR-SARShip-2.0}
\label{sec:air_sarship_results}

表~\ref{tab:air_sarship_comparison}
列出 AIR-SARShip-2.0 上的已有检测结果。
除文献报告值外，
本文分别保留两组已经完成的实验：
第一组采用 640 输入的历史切片配置；
第二组采用 416 输入和 240 轮训练配置。
两组实验中的 RGSP 表现方向不同，
说明当前结论需要结合具体训练协议理解。

\Needspace{10\baselineskip}
\begingroup
\small
\setlength{\tabcolsep}{5pt}
\setlength{\LTcapwidth}{\linewidth}
\begin{longtable}{@{}lrrr@{}}
\caption{AIR-SARShip-2.0 目标检测结果（\%）。
本文两组已完成实验通过上标区分，
新增 640 输入、300 轮实验的结果尚未纳入。}\label{tab:air_sarship_comparison}\\
\toprule
Method & AP$_{50}$ & AP$_{75}$ & AP$_{50:95}$ \\
\midrule
\endfirsthead
\multicolumn{4}{c}{表~\thetable（续）}\\
\toprule
Method & AP$_{50}$ & AP$_{75}$ & AP$_{50:95}$ \\
\midrule
\endhead
\midrule
\multicolumn{4}{r}{接下页}\\
\endfoot
\bottomrule
\endlastfoot
ImageNet-SSD      & 83.00 & 46.22 & 46.49 \\
OSD-SSD           & 84.70 & 56.09 & 50.96 \\
OSM-SSD           & 84.49 & 55.09 & 51.14 \\
WBF-DM            & 85.43 & 61.56 & 54.78 \\
FMP-Net           & 83.00 & 51.90 & 48.50 \\
\midrule
Mamba-YOLO-T$^{a}$
                  & 47.590 & 7.463  & 17.375 \\
\textbf{RGSP (Ours)}$^{a}$
                  & 51.559 & 9.087  & 19.291 \\
Mamba-YOLO-T$^{b}$
                  & 47.156 & 19.403 & 22.947 \\
\textbf{RGSP (Ours)}$^{b}$
                  & 36.423 & 15.284 & 18.172 \\
\end{longtable}
\noindent
\smallskip
\begin{minipage}{\linewidth}
\footnotesize
\textbf{表~\ref{tab:air_sarship_comparison} 注：}
$^{a}$历史 640 输入配置：
原图按 210/45/45 划分训练、验证与测试集，
采用 640 像素切片、512 像素步长，
从头训练并依据验证集选择权重。
$^{b}$416 输入配置：
原图按 210/90 划分训练与测试集，
采用 512 像素切片、256 像素步长，
从头训练 240 轮，
使用最后一轮权重进行测试。
两组配置均保留空背景切片。
外部文献与本文的完整协议尚未对齐，
不进行跨协议最优值标注。
\end{minipage}
\endgroup
\FloatBarrier

在配置 $a$ 中，
RGSP 相对对应基线的 AP$_{50:95}$
提高 $1.916$ 个百分点；
在配置 $b$ 中则降低 $4.775$ 个百分点。
因此，已有实验尚不支持
RGSP 在该数据集上具有稳定的性能增益。
\FloatBarrier
\subsection{语义分割}
\label{sec:segmentation_results}

在 OpenEarthMap-SAR 数据集上评估 RGSP 的语义分割性能。
本文以 VMamba-Tiny 为骨干网络，
采用 UPerNet-lite 解码器，
比较原始网络与接入 RGSP 后的分割结果。
评价指标包括平均交并比（mIoU）、
频率加权交并比（FWIoU）
和总体准确率（OA）。

表~\ref{tab:openearthmap_sar_comparison}
汇总外部方法及本文配置的结果。
外部比较包括 VMamba、U-Net + EfficientNet-B4
和 SegFormer，
其数值采用 OpenEarthMap-SAR 基准论文报告的
SAR 单模态、P+R5 设置，
即使用伪标签并加入每个区域 5 张真实标注图像。
本文分别列出早期批次和有效批量为 64 的配对实验，
各批次均保留对应的基础网络结果。

\Needspace{10\baselineskip}
\begingroup
\small
\setlength{\tabcolsep}{4pt}
\setlength{\LTcapwidth}{\linewidth}
\begin{longtable}{@{}llrrr@{}}
\caption{OpenEarthMap-SAR 语义分割结果（\%）。
方法后注明原方法的发表出处与年份；
外部数值来自 OpenEarthMap-SAR 基准论文。
粗体仅标识本文方法，不表示最优结果。}\label{tab:openearthmap_sar_comparison}\\
\toprule
Method & Setting & mIoU & FWIoU & OA \\
\midrule
\endfirsthead
\multicolumn{5}{c}{表~\thetable（续）}\\
\toprule
Method & Setting & mIoU & FWIoU & OA \\
\midrule
\endhead
\midrule
\multicolumn{5}{r}{接下页}\\
\endfoot
\bottomrule
\endlastfoot
VMamba (NeurIPS 2024)
& SAR, P+R5
& 35.08 & -- & -- \\

U-Net + EfficientNet-B4 (Benchmark 2025)
& SAR, P+R5
& 36.84 & -- & -- \\

SegFormer (NeurIPS 2021)
& SAR, P+R5
& 37.23 & -- & -- \\

\midrule
VMamba-Tiny + UPerNet-lite
& Configuration A
& 36.401 & 37.075 & 53.553 \\

\textbf{VMamba-Tiny + UPerNet-lite + RGSP}
& Configuration A
& 36.532 & 36.470 & 52.683 \\

VMamba-Tiny + UPerNet-lite
& Configuration B
& 36.551 & 36.120 & 52.316 \\

\textbf{VMamba-Tiny + UPerNet-lite + RGSP}
& Configuration B
& 35.898 & 35.220 & 51.764 \\
\end{longtable}
\noindent
\smallskip
\begin{minipage}{\linewidth}
\footnotesize
\textbf{表~\ref{tab:openearthmap_sar_comparison} 注：}
Configuration A 为已有早期配对实验，
RGSP 未采用额外模块预热，
并依据验证集 mIoU 选择权重。
Configuration B 为有效批量 64 的配对实验，
采用单次批量 32、梯度累积 2 次，
同样不采用额外模块预热。
本文结果为单种子记录。
外部方法与本文配置的标注设置及数据划分
尚未完成逐项对齐，因此不进行跨协议最优值标注。
``--'' 表示未报告。
\end{minipage}
\endgroup
\FloatBarrier

在配置 A 中，
RGSP 的 mIoU 从 $36.401\%$ 提高至 $36.532\%$，
增加 $0.131$ 个百分点，
但 FWIoU 和 OA 分别下降
$0.605$ 和 $0.870$ 个百分点。
在配置 B 中，
RGSP 的 mIoU 为 $35.898\%$，
相较对应基线下降 $0.653$ 个百分点，
FWIoU 和 OA 也有所下降。
因此，现有分割实验仅在早期配置的 mIoU 指标上
观察到小幅改善，
尚未显示跨配置一致的性能增益。
\FloatBarrier
\subsection{变化检测}
\label{sec:change_detection_results}

在 Ottawa 和 Yellow River 数据集上评估 RGSP
在 SAR 图像变化检测中的表现。
比较方法包括 DDNet、LANTNet、CAMixer、
WBANet、MDDNet 和 SARCDNet。
外部方法的 F1 数值取自 SARCDNet 论文的比较表，
方法名称后标注其发表来源和年份。
本文实验报告 F1、IoU、总体准确率（OA）、
Precision 和 Recall。

为考察不同配置下的表现，
分别列出批大小对照和早期状态维数对照。
同一重跑批次内，RGSP 与对应批大小的基线进行比较；
早期归档结果单独标注。
外部方法与本文实验的训练和评估协议尚未完全对齐，
因此文献结果用于参考，不据此作同协议性能排名。

\subsubsection{Ottawa}
\label{sec:ottawa_results}

表~\ref{tab:ottawa_comparison}
列出 Ottawa 上的变化检测结果。
在批大小分别为 16、32、64 和 96 的重跑实验中，
RGSP 的 F1 分别为
92.249\%、91.827\%、91.981\% 和 92.039\%，
相较各自基线分别提高
0.158、0.713、0.897 和 0.630 个百分点。
其中，BS16 配置的 Recall
由 94.000\% 提高至 96.003\%，
Precision 则由 90.259\% 降至 88.777\%。

早期归档批次呈现不同趋势：
状态维数为 32 和 16 的 RGSP 配置
均低于该批次基线。
因此，现有结果表明 RGSP 在重跑批次中获得小幅改善，
但这一趋势尚未在所有配置中保持一致。

\Needspace{10\baselineskip}
\begingroup
\scriptsize
\setlength{\tabcolsep}{2pt}
\renewcommand{\arraystretch}{1.12}
\setlength{\LTcapwidth}{\linewidth}
\begin{xltabular}{\linewidth}{@{}Xlrrrrr@{}}
\caption{Ottawa 变化检测结果（\%）。
所有指标均为越高越好。
粗体仅标识本文方法，不表示最优结果。}\label{tab:ottawa_comparison}\\
\toprule
Method & Setting & F1 & IoU & OA & Precision & Recall \\
\midrule
\endfirsthead
\multicolumn{7}{c}{表~\thetable（续）}\\
\toprule
Method & Setting & F1 & IoU & OA & Precision & Recall \\
\midrule
\endhead
\midrule
\multicolumn{7}{r}{接下页}\\
\endfoot
\bottomrule
\endlastfoot
DDNet (GRSL, 2021 online)
& Literature & 94.1 & -- & -- & -- & -- \\
LANTNet (GRSL, 2022)
& Literature & 94.5 & -- & -- & -- & -- \\
CAMixer (GRSL, 2023)
& Literature & 93.9 & -- & -- & -- & -- \\
WBANet (GRSL, 2024)
& Literature & 94.2 & -- & -- & -- & -- \\
MDDNet (ICASSP, 2025)
& Literature & 94.5 & -- & -- & -- & -- \\
SARCDNet (Scientific Reports, 2026)
& Literature & 94.9 & -- & -- & -- & -- \\
\midrule
Baseline
& A: BS16, d32
& 92.758 & 86.495 & 98.491 & 90.990 & 94.597 \\
\textbf{Baseline + RGSP}
& A: BS16, d32
& 92.110 & 85.374 & 98.350 & 90.061 & 94.254 \\
\textbf{Baseline + RGSP}
& A: BS16, d16
& 92.173 & 85.483 & 98.346 & 89.238 & 95.308 \\
\midrule
Baseline
& B: BS16
& 92.091 & 85.342 & 98.351 & 90.259 & 94.000 \\
\textbf{Baseline + RGSP}
& B: BS16, d16
& 92.249 & 85.612 & 98.352 & 88.777 & 96.003 \\
Baseline
& B: BS32
& 91.114 & 83.678 & 98.102 & 87.319 & 95.254 \\
\textbf{Baseline + RGSP}
& B: BS32, d16
& 91.827 & 84.889 & 98.272 & 88.818 & 95.047 \\
Baseline
& B: BS64
& 91.084 & 83.628 & 98.096 & 87.323 & 95.184 \\
\textbf{Baseline + RGSP}
& B: BS64, d16
& 91.981 & 85.152 & 98.320 & 89.764 & 94.310 \\
Baseline
& B: BS96
& 91.409 & 84.177 & 98.173 & 87.953 & 95.147 \\
\textbf{Baseline + RGSP}
& B: BS96, d16
& 92.039 & 85.251 & 98.333 & 89.841 & 94.346 \\
\end{xltabular}
\noindent
\smallskip
\begin{minipage}{\linewidth}
\footnotesize
\textbf{表~\ref{tab:ottawa_comparison} 注：}
A 表示早期归档批次，B 表示批大小重跑批次；
BS 为 batch size，d 为状态维数。
重跑批次结果为 seed 0 的单次记录。
GRSL 指 IEEE Geoscience and Remote Sensing Letters。
外部 F1 数值来自 SARCDNet 的 Table 2\cite{sarcdnet2026}；
``--'' 表示当前汇总中未收录该指标。
不同批次分别比较，不跨批次计算模块增益。
\end{minipage}
\endgroup
\FloatBarrier

\subsubsection{Yellow River}
\label{sec:yellow_river_results}

表~\ref{tab:yellow_river_comparison}
列出 Yellow River 上的变化检测结果。
早期归档批次中，
状态维数为 32 的 RGSP 配置取得
87.689\% 的 F1，
对应基线为 88.079\%，
下降 0.390 个百分点；
状态维数为 16 时，F1 为 86.308\%。

在 BS16、BS32、BS64 和 BS96 的重跑实验中，
RGSP 的 F1 相较对应基线分别下降
1.477、2.719、3.642 和 3.543 个百分点。
其中 BS16、BS96 的 Precision 提高；BS32、BS64 分别由 91.293\%、91.834\% 降至 90.000\%、91.671\%，
但 Recall 均有所下降，
最终 F1 未获得改善。
因此，当前实验未显示 RGSP 在
Yellow River 上的性能增益。

\Needspace{10\baselineskip}
\begingroup
\scriptsize
\setlength{\tabcolsep}{2pt}
\renewcommand{\arraystretch}{1.12}
\setlength{\LTcapwidth}{\linewidth}
\begin{xltabular}{\linewidth}{@{}Xlrrrrr@{}}
\caption{Yellow River 变化检测结果（\%）。
所有指标均为越高越好。
粗体仅标识本文方法，不表示最优结果。}\label{tab:yellow_river_comparison}\\
\toprule
Method & Setting & F1 & IoU & OA & Precision & Recall \\
\midrule
\endfirsthead
\multicolumn{7}{c}{表~\thetable（续）}\\
\toprule
Method & Setting & F1 & IoU & OA & Precision & Recall \\
\midrule
\endhead
\midrule
\multicolumn{7}{r}{接下页}\\
\endfoot
\bottomrule
\endlastfoot
DDNet (GRSL, 2021 online)
& Literature & 87.3 & -- & -- & -- & -- \\
LANTNet (GRSL, 2022)
& Literature & 85.8 & -- & -- & -- & -- \\
CAMixer (GRSL, 2023)
& Literature & 85.2 & -- & -- & -- & -- \\
WBANet (GRSL, 2024)
& Literature & 87.6 & -- & -- & -- & -- \\
MDDNet (ICASSP, 2025)
& Literature & 85.8 & -- & -- & -- & -- \\
SARCDNet (Scientific Reports, 2026)
& Literature & 90.0 & -- & -- & -- & -- \\
\midrule
Baseline
& A: BS16, d32
& 88.079 & 78.697 & 97.758 & 90.666 & 85.635 \\
\textbf{Baseline + RGSP}
& A: BS16, d32
& 87.689 & 78.077 & 97.683 & 90.227 & 85.290 \\
\textbf{Baseline + RGSP}
& A: BS16, d16
& 86.308 & 75.914 & 97.478 & 90.885 & 82.170 \\
\midrule
Baseline
& B: BS16
& 88.368 & 79.160 & 97.817 & 91.205 & 85.702 \\
\textbf{Baseline + RGSP}
& B: BS16, d16
& 86.891 & 76.820 & 97.584 & 91.422 & 82.787 \\
Baseline
& B: BS32
& 86.273 & 75.859 & 97.483 & 91.293 & 81.776 \\
\textbf{Baseline + RGSP}
& B: BS32, d16
& 83.554 & 71.753 & 97.031 & 90.000 & 77.971 \\
Baseline
& B: BS64
& 86.968 & 76.941 & 97.606 & 91.834 & 82.591 \\
\textbf{Baseline + RGSP}
& B: BS64, d16
& 83.326 & 71.419 & 97.043 & 91.671 & 76.374 \\
Baseline
& B: BS96
& 83.484 & 71.650 & 97.002 & 89.377 & 78.320 \\
\textbf{Baseline + RGSP}
& B: BS96, d16
& 79.941 & 66.585 & 96.546 & 91.218 & 71.146 \\
\end{xltabular}
\noindent
\smallskip
\begin{minipage}{\linewidth}
\footnotesize
\textbf{表~\ref{tab:yellow_river_comparison} 注：}
A 表示早期归档批次，B 表示批大小重跑批次；
BS 为 batch size，d 为状态维数。
重跑批次结果为 seed 0 的单次记录。
外部 F1 数值来自 SARCDNet 的 Table 2\cite{sarcdnet2026}；
``--'' 表示当前汇总中未收录该指标。
文献结果与本文实验的协议尚未完全对齐，
不据此作同协议性能排名。
\end{minipage}
\endgroup
\FloatBarrier

\FloatBarrier
\section{后续验收项}
完成新 SSDD 36-run 的逐项训练/测试与关闭条件核验，补齐公平第二/第三数据集和外部轻量基线，统一效率协议，并检查真实预测和引导信号。未知结果保持未知，不根据中间结果修改控制定义。需要继续核验的文献行与跨任务配置详见来源登记，不把本文编辑等同于实验实施。
\begin{thebibliography}{99}
\interlinepenalty=10000
\bibitem{sarcdnet2026}
V. D. Kevala, V. Mukundan, S. Nedungatt, S. Suresh, and S. Lal, SARCDNet-an enhanced deep learning network for change detection from bi-temporal SAR images, \emph{Scientific Reports}, vol. 16, art. 1826, 2026 (online Dec. 31, 2025). doi:10.1038/s41598-025-31488-y.
\end{thebibliography}

\end{document}
